\documentclass[11pt]{article}
\usepackage[letterpaper,margin=0.85in]{geometry}\usepackage[T1]{fontenc}\usepackage[utf8]{inputenc}\usepackage{amsmath,amssymb,enumitem,booktabs,array,xcolor,fancyhdr,tabularx}
\setlist{itemsep=2pt,topsep=3pt}\pagestyle{fancy}\fancyhf{}\renewcommand{\headrulewidth}{0.3pt}
\fancyhead[L]{\small Counting Particles (template)}\fancyhead[R]{\small Make-up key}\fancyfoot[C]{\small\thepage}
\definecolor{ans}{RGB}{28,125,74}\newcommand{\A}[1]{\textcolor{ans}{\textbf{#1}}}
\begin{document}

\section*{Counting Particles (template)}
\textbf{You missed this lesson. Copy this into your notes. That is all you need to do.}\par\medskip
\subsection*{Question of the day}
How do we count things too small to see?\par
\subsection*{Targets}\begin{itemize}\item I can say what a mole is.
\item I can turn grams into moles with molar mass.
\item I can turn moles into a count of particles.\end{itemize}
\subsection*{Concepts}\begin{itemize}
\item \textbf{Mole} (1 mol = 6.02 $\times$ 10\textsuperscript{2}\textsuperscript{3}): A counting unit, like a dozen. One mole is 6.02 $\times$ 10\textsuperscript{2}\textsuperscript{3} particles.
\item \textbf{Molar mass} (n = m / M): Grams in one mole. Add the atomic masses in the formula.
\end{itemize}
\subsection*{What the class predicted, and what is true}\begin{itemize}
\item One mole of water and one mole of iron. Which holds more particles? \A{The same number.} A mole is a count. One mole of anything is 6.02 $\times$ 10\textsuperscript{2}\textsuperscript{3} particles.
\item Same two samples: one mole of water and one mole of iron. Which has more mass? \A{Iron.} Same count, but each iron atom is heavier than a water molecule (56 g/mol vs 18 g/mol).
\end{itemize}
\section*{Huddle notes (answers filled in)}
\subsection*{Definitions}\begin{itemize}
\item \textbf{Mole:} The amount that holds \A{6.02 $\times$ 10\textsuperscript{2}\textsuperscript{3}} particles.
\item \textbf{Molar mass:} The \A{mass} of one mole, in grams per mole.
\end{itemize}
\subsection*{Worked example}
How many moles are in 36 g of water? (M = 18 g/mol)\\
\begin{itemize}
\item n = m / M
\item n = 36 / 18 = \A{2} mol
\item Particles: 2 $\times$ 6.02 $\times$ 10\textsuperscript{2}\textsuperscript{3} = \A{12.04} $\times$ 10\textsuperscript{2}\textsuperscript{3}
\end{itemize}
\subsection*{Label the water molecule}
\begin{enumerate}
\item The \A{oxygen} atom.
\item A \A{hydrogen} atom.
\item The \A{bond} that holds them.
\end{enumerate}
\section*{One worked sample of each skill}
\subsection*{Grams to moles}
\noindent\textbf{Level 1.} A sample of water has a mass of 54 g. Water has a molar mass of 18 g/mol. How many moles is that?\\ \A{Answer: 3 mol} \quad n = m / M = 54 / 18 = 3 mol.\par\medskip
\noindent\textbf{Level 2.} A sample of carbon has a mass of 36 g. Its molar mass is 12 g/mol. How many moles is that?\\ \A{Answer: 3 mol} \quad n = 36 / 12 = 3 mol.\par\medskip
\noindent\textbf{Level 3.} A sample of CH? has a mass of 4 g. The molar mass of CH? is 16 g/mol. How many moles is the sample?\\ \A{Answer: 0.25 mol} \quad n = 4 / 16 = 0.25 mol.\par\medskip
\noindent\textbf{Level 4.} You weigh out 116 g of a compound with M = 58 g/mol. How many moles do you have? Answer in mol.\\ \A{Answer: 2 mol} \quad n = 116 / 58 = 2 mol.\par\medskip
\subsection*{Moles to particles}
\noindent\textbf{Level 1.} How many particles are in 2 mol? Give the answer in units of 10\textsuperscript{2}\textsuperscript{3} (so 6.02 $\times$ 10\textsuperscript{2}\textsuperscript{3} is written 6.02).\\ \A{Answer: 12.04 $\times$ 10\textsuperscript{2}\textsuperscript{3}} \quad 2 $\times$ 6.02 = 12.04, so 12.04 $\times$ 10\textsuperscript{2}\textsuperscript{3} particles.\par\medskip
\noindent\textbf{Level 2.} How many particles are in 0.5 mol? Answer in units of 10\textsuperscript{2}\textsuperscript{3}.\\ \A{Answer: 3.01 $\times$ 10\textsuperscript{2}\textsuperscript{3}} \quad 0.5 $\times$ 6.02 = 3.01 $\times$ 10\textsuperscript{2}\textsuperscript{3}.\par\medskip
\noindent\textbf{Level 3.} How many particles are in 4 mol? Answer in units of 10\textsuperscript{2}\textsuperscript{3}.\\ \A{Answer: 24.08 $\times$ 10\textsuperscript{2}\textsuperscript{3}} \quad 4 $\times$ 6.02 = 24.08 $\times$ 10\textsuperscript{2}\textsuperscript{3}.\par\medskip
\noindent\textbf{Level 4.} How many water molecules are in 54 g of water (M = 18 g/mol)? Answer in units of 10\textsuperscript{2}\textsuperscript{3}.\\ \A{Answer: 18.06 $\times$ 10\textsuperscript{2}\textsuperscript{3}} \quad n = 54 / 18 = 3 mol. Particles = 3 $\times$ 6.02 = 18.06 $\times$ 10\textsuperscript{2}\textsuperscript{3}.\par\medskip
\end{document}